\documentclass{standalone}
\usepackage{tikz}
\usetikzlibrary{intersections, calc, angles, quotes}

\begin{document}

\begin{tikzpicture}[scale=1.2]
    % Define coordinates
    \coordinate (O) at (0,0);
    \def\R{3.5}
    \coordinate (A) at (90:\R);
    \coordinate (B) at (225:\R);
    \coordinate (C) at (315:\R);
    \coordinate (E) at (45:\R);
    
    % H is the projection of A on BC. Since B and C are symmetric wrt y-axis, H is on y-axis.
    % y-coordinate of B and C is -R*sin(45)
    \coordinate (H) at (0, {-\R*sin(45)});
    
    % Line AC
    \path[name path=AC] (A) -- (C);
    % Line BE (Diameter)
    \path[name path=BE] (B) -- (E);
    
    % Intersection D
    \path[name intersections={of=AC and BE, by=D}];
    
    % F is the foot of the perpendicular from B to AC
    \path[name path=BF_line] (B) -- ($(B)!10!(A)$); % Dummy path to define direction, actually we need perpendicular
    % Let's calculate F properly using projection
    \coordinate (F) at ($(A)!(B)!(C)$);
    
    % Draw Circle
    \draw (O) circle (\R);
    
    % Draw Triangle ABC
    \draw (A) -- (B) -- (C) -- cycle;
    
    % Draw Altitude AH
    % AO is dashed, OH is solid based on visual inspection
    \draw[dashed] (A) -- (O);
    \draw (O) -- (H);
    
    % Draw Line BE with gap DO
    % The instruction says: "Delete the sub-segment DO on the continuous straight line."
    % "The line segments outside DO on that straight line must be preserved."
    % So we draw B -- O and D -- E.
    \draw (B) -- (O);
    \draw (D) -- (E);
    
    % Draw Line OC
    \draw (O) -- (C);
    
    % Draw Line BF (perpendicular to AC)
    \draw (B) -- (F);
    
    % Right angle mark at F
    \draw ($(F)!3mm!(A)$) -- ($(F)!3mm!(A)!3mm!-90:(A)$) -- ($(F)!3mm!(B)$);
    
    % Labels
    \node[above] at (A) {$A$};
    \node[below left] at (B) {$B$};
    \node[below right] at (C) {$C$};
    \node[above right] at (D) {$D$};
    \node[above right] at (E) {$E$};
    \node[below right] at (F) {$F$};
    \node[below] at (H) {$H$};
    \node[above left] at (O) {$O$};
    
\end{tikzpicture}

\end{document}
